\chapter{Conclusiones}
\label{chap:conclusions}

\lettrine{E}{n} este último capítulo se presenta una reflexión final sobre las prácticas realizadas, incluyendo la relación con los conocimientos adquiridos durante el grado y una valoración personal de la experiencia.

\section{Relación con la titulación}
\label{sec:titulacion}

Durante el desarrollo de las prácticas, varias asignaturas del grado han resultado especialmente útiles:

\begin{itemize}
    \item \textbf{Aprendizaje Automático}: Los conocimientos adquiridos en esta asignatura fueron fundamentales para entender y trabajar con los modelos de ML, especialmente en la parte de evaluación de métricas y entrenamiento. Las herramientas como MLflow habrían sido muy útiles para los proyectos de la asignatura.
    
    \item \textbf{Concurrencia y Paralelismo}: Los conceptos aprendidos fueron esenciales para entender y depurar problemas relacionados con threads y procesos, especialmente en el desarrollo de interfaces y en la gestión de procesos de entrenamiento.
    
    \item \textbf{Paradigmas de Programación}: La experiencia con diferentes paradigmas ayudó en la comprensión y desarrollo de código Python, especialmente en aspectos como type hints y programación funcional.
    
    \item \textbf{Visión Artificial}: Los fundamentos teóricos fueron valiosos para comprender los sistemas de visión artificial y reconocimiento de patrones.
    
    \item \textbf{Interfaces Persona Máquina}: Los conocimientos sobre diseño de interfaces facilitaron la implementación de mejoras en las herramientas de configuración.
    
    \item \textbf{Algoritmos}: La base algorítmica fue importante para optimizar procesos y mejorar la eficiencia del código.
\end{itemize}

\section{Valoración de la práctica}
\label{sec:valoracion}

Las prácticas han sido una experiencia muy enriquecedora que ha permitido:

\begin{enumerate}
    \item \textbf{Experiencia profesional real}:
    \begin{itemize}
        \item Trabajo en equipo usando metodologías ágiles
        \item Gestión de proyectos con JIRA y GitLab
        \item Participación en code reviews y merge requests
    \end{itemize}
    
    \item \textbf{Aprendizaje técnico}:
    \begin{itemize}
        \item Profundización en Python y buenas prácticas de desarrollo
        \item Uso de herramientas profesionales (Docker, Git, MLflow, DVC)
        \item Experiencia en ML/CV y procesamiento de audio
    \end{itemize}
    
    \item \textbf{Desarrollo personal}:
    \begin{itemize}
        \item Mejora en la capacidad de resolución de problemas
        \item Desarrollo de autonomía y responsabilidad
        \item Comunicación efectiva en equipo
    \end{itemize}
\end{enumerate}

\section{Lecciones aprendidas}
\label{sec:lecciones}

Entre las principales lecciones aprendidas destacan:

\begin{itemize}
    \item \textbf{Gestión de dependencias}: La importancia de mantener un control estricto de las versiones de las dependencias y el uso de entornos virtuales.
    
    \item \textbf{Documentación}: La necesidad de mantener una documentación clara y actualizada, especialmente en proyectos colaborativos.
    
    \item \textbf{Control de versiones}: La importancia de una buena gestión de ramas y commits en Git para mantener un desarrollo ordenado.
    
    \item \textbf{Optimización de recursos}: La gestión eficiente de recursos computacionales, especialmente en entrenamientos de modelos ML.
    
    \item \textbf{Testing}: La importancia de las pruebas exhaustivas antes de pasar código a producción.
\end{itemize}